% ============================================================
%  CHAPITRE — ANALYSE DES BESOINS ET CONCEPTION
% ============================================================
\chapter{Analyse des Besoins et Conception de la Solution}

\section*{Introduction}

Ce chapitre présente la démarche de conception qui a guidé le développement
de la plateforme GRC AI Analyzer. Après avoir établi le contexte métier et
la problématique dans les chapitres précédents, il s'agit ici de traduire
les exigences opérationnelles en artefacts de conception formels : analyse
des besoins fonctionnels et non fonctionnels, diagramme des cas d'utilisation,
diagramme de classes, diagramme de séquence et diagramme d'activité.
Ces modèles constituent la référence conceptuelle de toute la phase
d'implémentation.

% ============================================================
\section{Analyse des Besoins}
% ============================================================

\subsection{Identification des acteurs}

La plateforme GRC AI Analyzer implique trois acteurs principaux dont les
périmètres d'action sont distincts :

\begin{itemize}
  \item \textbf{L'Analyste GRC} — acteur central du système. Il uploade
    les rapports (PDF, DOCX, XLSX), déclenche les analyses, consulte les
    résultats (matrice des risques, taux de conformité, recommandations),
    interagit avec l'assistant RAG et exporte les livrables.

  \item \textbf{L'Encadrant / Manager} — profil à accès en lecture. Il
    consulte le tableau de bord, visualise le score de maturité GRC global,
    suit l'avancement des recommandations et accède aux rapports d'analyse
    générés par les analystes.

  \item \textbf{L'Administrateur} — responsable de la gouvernance applicative.
    Il gère les comptes utilisateurs (création, attribution des rôles
    \texttt{ADMIN} / \texttt{ANALYSTE}), consulte la piste d'audit et
    peut configurer les référentiels de conformité.

  \item \textbf{Mistral AI} — acteur externe (API tierce). Le modèle
    \texttt{mistral-large-latest} est sollicité pour l'extraction des
    risques, les vérifications de conformité, la génération des
    recommandations et la production du résumé exécutif.
\end{itemize}

\subsection{Besoins fonctionnels}

Les besoins fonctionnels ont été organisés en cinq modules cohérents :

\subsubsection*{Module 1 — Gestion des rapports}

\begin{itemize}
  \item \textbf{BF-01} : L'analyste doit pouvoir uploader un fichier au
    format PDF, DOCX ou XLSX dont la taille ne dépasse pas la limite
    configurée. Le système valide l'extension et rejette tout format non
    supporté.
  \item \textbf{BF-02} : Le système enregistre les métadonnées du rapport
    (nom, format, date d'upload, statut) dans la base de données et stocke
    le fichier binaire sur le serveur de fichiers local.
  \item \textbf{BF-03} : L'analyste peut consulter l'historique de ses
    rapports et filtrer par statut (\textit{en\_attente}, \textit{en\_cours},
    \textit{termine}, \textit{erreur}).
  \item \textbf{BF-04} : L'analyste peut supprimer un rapport, ce qui
    entraîne la suppression en cascade des analyses associées, des fichiers
    physiques et des collections vectorielles ChromaDB correspondantes.
\end{itemize}

\subsubsection*{Module 2 — Analyse IA automatisée}

\begin{itemize}
  \item \textbf{BF-05} : À partir d'un rapport uploadé, l'analyste peut
    lancer une analyse automatique. L'API répond immédiatement (HTTP 202)
    et l'analyse s'exécute en tâche de fond.
  \item \textbf{BF-06} : Le système extrait le texte brut du document
    (page par page pour PDF via PyMuPDF, paragraphe par paragraphe pour
    DOCX via python-docx, feuille par feuille pour XLSX via openpyxl).
  \item \textbf{BF-07} : Le système découpe le texte en fragments
    (\textit{chunks}) de 1 500 caractères avec un chevauchement de 200
    caractères, puis les indexe dans ChromaDB via le modèle d'embedding
    multilingue \texttt{paraphrase-multilingual-MiniLM-L12-v2}.
  \item \textbf{BF-08} : Le système extrait les risques de manière
    exhaustive : chaque alerte, vulnérabilité, écart ou finding du document
    constitue un risque distinct, caractérisé par un libellé, une
    description technique, une catégorie (\textit{CYBER, OPERATIONNEL,
    LEGAL, FINANCIER, RH}), une probabilité (1-5), un impact (1-5) et
    un score calculé (probabilité × impact).
  \item \textbf{BF-09} : La sévérité est déterminée automatiquement selon
    les seuils : score $\geq$ 20 → CRITIQUE ; $\geq$ 12 → ÉLEVÉ ;
    $\geq$ 6 → MOYEN ; $<$ 6 → FAIBLE.
  \item \textbf{BF-10} : Le système vérifie la conformité du document par
    rapport à trois référentiels en parallèle : ISO 27001:2013 (14 domaines
    de l'Annexe A), RGPD (12 domaines), Loi 09-08 marocaine (11 domaines).
    Chaque domaine reçoit un statut (CONFORME / PARTIEL / NON\_CONFORME)
    et un taux de conformité.
  \item \textbf{BF-11} : Le système génère une recommandation priorisée
    pour chaque risque identifié, incluant un type d'action
    (\textit{QUICK\_WIN} ou \textit{LONG\_TERME}) et un effort estimé
    (\textit{FAIBLE, MOYEN, ÉLEVÉ}).
  \item \textbf{BF-12} : Le système produit un résumé exécutif structuré
    (350-450 mots), destiné à un CISO, synthétisant les conclusions
    clés, les risques prioritaires, l'état de conformité et les actions
    immédiates recommandées.
  \item \textbf{BF-13} : Le score de maturité GRC global est calculé comme
    la moyenne arithmétique des taux de conformité des trois référentiels.
\end{itemize}

\subsubsection*{Module 3 — Consultation et visualisation}

\begin{itemize}
  \item \textbf{BF-14} : L'analyste peut consulter la matrice des risques
    (visualisation Probabilité × Impact) via un composant Recharts dédié.
  \item \textbf{BF-15} : L'analyste peut visualiser le taux de conformité
    par référentiel via un radar chart (composant \texttt{ComplianceRadar}).
  \item \textbf{BF-16} : L'analyste peut lire les recommandations classées
    par priorité (CRITIQUE → HAUTE → MOYENNE → FAIBLE) et mettre à jour
    leur statut de traitement (\textit{À\_FAIRE} → \textit{EN\_COURS}
    → \textit{CLÔTURÉ}).
  \item \textbf{BF-17} : Le tableau de bord affiche un score de maturité
    global, le nombre de risques par sévérité et les taux de conformité
    de la dernière analyse.
\end{itemize}

\subsubsection*{Module 4 — Assistant RAG conversationnel}

\begin{itemize}
  \item \textbf{BF-18} : L'analyste peut poser des questions en langage
    naturel (français ou anglais) sur le contenu du rapport. Le système
    récupère les 5 fragments les plus pertinents dans ChromaDB par
    similarité cosinus, construit un contexte et interroge Mistral pour
    générer une réponse sourcée.
  \item \textbf{BF-19} : Les réponses du chat citent explicitement les
    extraits du rapport sur lesquels elles s'appuient.
\end{itemize}

\subsubsection*{Module 5 — Export}

\begin{itemize}
  \item \textbf{BF-20} : L'analyste peut exporter l'analyse complète en
    PDF (rapport structuré avec page de couverture, résumé exécutif,
    tableau des risques, résultats de conformité et recommandations),
    généré via ReportLab.
  \item \textbf{BF-21} : L'analyste peut exporter le registre des risques
    en fichier Excel (.xlsx) avec quatre onglets (Résumé, Risques,
    Conformité, Recommandations), généré via XlsxWriter avec mise en
    forme conditionnelle par sévérité.
\end{itemize}

\subsubsection*{Module 6 — Authentification et administration}

\begin{itemize}
  \item \textbf{BF-22} : Le système implémente une authentification par
    JWT (JSON Web Token). Le token est signé avec l'algorithme HS256,
    stocké côté client dans le \textit{localStorage} via Zustand, et
    transmis dans l'en-tête \texttt{Authorization: Bearer} à chaque
    requête.
  \item \textbf{BF-23} : Les mots de passe sont hachés avec bcrypt avant
    stockage. La vérification à la connexion compare le hash stocké sans
    jamais manipuler le mot de passe en clair.
  \item \textbf{BF-24} : L'administrateur peut créer et gérer les
    utilisateurs, leur attribuer le rôle \texttt{ADMIN} ou
    \texttt{ANALYSTE}. Un analyste ne peut accéder qu'à ses propres
    rapports ; un administrateur a une vue globale.
\end{itemize}

\subsection{Besoins non fonctionnels}

\subsubsection*{Performance et scalabilité}

\begin{itemize}
  \item \textbf{BNF-01} : L'API doit répondre aux requêtes de consultation
    en moins de 500 ms en conditions normales. L'upload et le déclenchement
    de l'analyse retournent immédiatement (HTTP 202) sans bloquer le
    client.
  \item \textbf{BNF-02} : Le pipeline d'analyse exploite la concurrence
    via \texttt{asyncio.gather} : l'indexation ChromaDB, l'extraction
    des risques et les trois vérifications de conformité s'exécutent en
    parallèle (Phase 1) ; la génération des recommandations et du résumé
    exécutif s'exécutent elles aussi en parallèle (Phase 2).
  \item \textbf{BNF-03} : Les appels à l'API Mistral sont sérialisés par
    un sémaphore global (\texttt{asyncio.Semaphore(1)}) avec un délai
    minimum de 1,5 seconde entre appels consécutifs pour respecter les
    limites de débit (\textit{rate limits}). Un mécanisme de backoff
    exponentiel gère les erreurs HTTP 429 (délai de 30 s × numéro
    de tentative).
\end{itemize}

\subsubsection*{Fiabilité et résilience}

\begin{itemize}
  \item \textbf{BNF-04} : L'indexation RAG dans ChromaDB est
    \textit{non-bloquante} : un échec de ChromaDB n'interrompt pas
    l'analyse, la plateforme continue sans fonctionnalité de chat.
  \item \textbf{BNF-05} : En cas d'erreur dans le pipeline, le statut
    de l'analyse est mis à jour en \textit{erreur} dans la base de
    données et l'erreur est journalisée via Loguru, sans lever
    d'exception non gérée.
  \item \textbf{BNF-06} : Toutes les sessions de base de données du
    pipeline asynchrone s'exécutent dans un contexte transactionnel
    dédié (\texttt{AsyncSessionLocal}), distinct de la session HTTP,
    pour éviter tout conflit de concurrence.
\end{itemize}

\subsubsection*{Sécurité}

\begin{itemize}
  \item \textbf{BNF-07} : Toutes les routes API (sauf \texttt{/auth/login}
    et \texttt{/auth/register}) sont protégées par la dépendance
    \texttt{get\_current\_user} qui vérifie la validité du JWT.
  \item \textbf{BNF-08} : Les tokens JWT ont une durée de vie de 480
    minutes (8 heures). Un intercepteur Axios côté frontend invalide la
    session locale dès réception d'une réponse HTTP 401.
  \item \textbf{BNF-09} : La politique CORS est configurée explicitement
    pour n'autoriser que l'origine du frontend (\texttt{localhost:5173}
    en développement, URL configurable en production).
\end{itemize}

\subsubsection*{Maintenabilité et extensibilité}

\begin{itemize}
  \item \textbf{BNF-10} : L'architecture adopte une séparation stricte
    des responsabilités : les routes FastAPI sont de simples façades,
    toute la logique métier réside dans la couche \texttt{services/},
    et la logique IA dans la couche \texttt{ai/}.
  \item \textbf{BNF-11} : Les migrations de base de données sont gérées
    par Alembic, permettant des évolutions du schéma sans perte de données.
  \item \textbf{BNF-12} : Les prompts Mistral sont externalisés dans des
    modules dédiés (\texttt{risk\_prompts.py}, \texttt{compliance\_prompts.py},
    \texttt{recommendation\_prompts.py}, \texttt{summary\_prompts.py}),
    facilitant leur maintenance et leur amélioration indépendante.
\end{itemize}

\subsubsection*{Internationalisation}

\begin{itemize}
  \item \textbf{BNF-13} : Le modèle d'embedding
    \texttt{paraphrase-multilingual-MiniLM-L12-v2} supporte 50+ langues,
    permettant l'analyse de documents en français, anglais ou arabe.
    L'assistant RAG répond dans la langue de la question posée.
\end{itemize}

% ============================================================
\section{Diagramme des Cas d'Utilisation}
% ============================================================

\subsection{Description générale}

Le diagramme des cas d'utilisation (figure~\ref{fig:use_case}) représente
l'ensemble des interactions fonctionnelles entre les acteurs et le système
GRC AI Analyzer. Il s'organise en six packages fonctionnels qui correspondent
aux modules identifiés lors de l'analyse des besoins.

\begin{figure}[htbp]
  \centering
  \includegraphics[width=0.92\textwidth]{figures/use_case_diagram.png}
  \caption{Diagramme des cas d'utilisation — GRC AI Analyzer}
  \label{fig:use_case}
\end{figure}

\subsection{Description des cas d'utilisation principaux}

\subsubsection*{UC-01 : Uploader un rapport}

\begin{tcolorbox}[colback=lightgray!20, colframe=darknavy, title=UC-01 — Uploader un rapport]
\textbf{Acteur principal :} Analyste GRC \\
\textbf{Précondition :} L'analyste est authentifié. \\
\textbf{Scénario nominal :}
\begin{enumerate}
  \item L'analyste sélectionne un fichier (PDF, DOCX ou XLSX) via l'interface.
  \item Le système valide le format (\texttt{validate\_file\_format}).
  \item Le fichier est sauvegardé sur le serveur sous
        \texttt{uploads/reports/\{uuid\}/\{nom\}}.
  \item Un enregistrement \texttt{Rapport} est créé en base avec le statut
        \textit{en\_attente}.
  \item L'API retourne l'identifiant du rapport (HTTP 201).
\end{enumerate}
\textbf{Exception :} Format non supporté → HTTP 400 avec message d'erreur.
\end{tcolorbox}

\subsubsection*{UC-02 : Lancer l'analyse automatique}

\begin{tcolorbox}[colback=lightgray!20, colframe=darknavy, title=UC-02 — Lancer l'analyse automatique]
\textbf{Acteur principal :} Analyste GRC \\
\textbf{Acteur secondaire :} Mistral AI \\
\textbf{Précondition :} Un rapport au statut \textit{en\_attente} existe. \\
\textbf{Scénario nominal :}
\begin{enumerate}
  \item L'analyste déclenche l'analyse via
        \texttt{POST /api/v1/analyses/lancer/\{rapport\_id\}}.
  \item Le système crée un enregistrement \texttt{Analyse} (statut
        \textit{en\_cours}) et retourne immédiatement HTTP 202.
  \item En tâche de fond : extraction du texte, chunking, indexation
        ChromaDB, extraction des risques et vérifications de conformité
        (Phase 1 parallèle), puis recommandations et résumé (Phase 2
        parallèle).
  \item L'analyse est marquée \textit{terminé} à la fin du pipeline.
\end{enumerate}
\textbf{Cas inclus :} Extraire les risques identifiés ; Scorer les risques ;
Vérifier la conformité ; Analyser les écarts (Gap Analysis) ;
Générer le résumé exécutif.
\end{tcolorbox}

\subsubsection*{UC-03 : Poser des questions via le Chat IA}

\begin{tcolorbox}[colback=lightgray!20, colframe=darknavy, title=UC-03 — Poser des questions sur le rapport (RAG Chat)]
\textbf{Acteur principal :} Analyste GRC \\
\textbf{Acteur secondaire :} Mistral AI \\
\textbf{Précondition :} L'analyse est terminée et les chunks ont été indexés. \\
\textbf{Scénario nominal :}
\begin{enumerate}
  \item L'analyste saisit une question dans l'onglet Assistant RAG.
  \item Le système encode la question via le modèle sentence-transformer.
  \item ChromaDB retourne les 5 chunks les plus proches par similarité cosinus.
  \item Le système construit le prompt contextuel et interroge Mistral.
  \item La réponse est affichée avec les extraits sources cités.
\end{enumerate}
\textbf{Cas inclus :} Obtenir des réponses citant les sources.
\end{tcolorbox}

\subsubsection*{UC-04 : Exporter les livrables}

L'analyste peut exporter deux types de livrables depuis l'interface :
\begin{itemize}
  \item \textbf{Export PDF} — rapport d'analyse complet généré par
    ReportLab, comprenant une page de couverture avec score de maturité,
    un résumé exécutif, un tableau des risques avec badges de sévérité
    colorés, les résultats de conformité par référentiel et le plan de
    recommandations.
  \item \textbf{Export Excel} — registre des risques au format XLSX
    (XlsxWriter) avec quatre onglets structurés et mise en forme
    conditionnelle par sévérité et priorité.
\end{itemize}

\subsubsection*{UC-05 : Suivre l'avancement des recommandations}

L'analyste peut faire progresser le statut de chaque recommandation
via l'interface (bouton de cycle d'état) :
\textit{À\_FAIRE} $\rightarrow$ \textit{EN\_COURS} $\rightarrow$
\textit{CLÔTURÉ}, en appelant
\texttt{PATCH /api/v1/recommandations/\{id\}/statut}.

% ============================================================
\section{Diagramme de Classes}
% ============================================================

\subsection{Vue d'ensemble}

Le diagramme de classes (figure~\ref{fig:class_diagram}) modélise les
entités persistantes du domaine et leurs relations. Il reflète fidèlement
le schéma de base de données PostgreSQL géré par SQLAlchemy (ORM asynchrone)
et versionné via Alembic.

\begin{figure}[htbp]
  \centering
  \includegraphics[width=0.95\textwidth]{figures/class_diagram.png}
  \caption{Diagramme de classes — modèle du domaine GRC AI Analyzer}
  \label{fig:class_diagram}
\end{figure}

\subsection{Description des classes principales}

\subsubsection*{Utilisateur}

La classe \texttt{Utilisateur} représente les comptes de la plateforme.
Chaque utilisateur possède un identifiant UUID, un nom, un email unique
(indexé), un mot de passe haché (bcrypt), un rôle énuméré
(\texttt{ADMIN} ou \texttt{ANALYSTE}) et une date de création.
Un utilisateur peut uploader plusieurs rapports et déclencher plusieurs
analyses (relations \textit{one-to-many}).

\subsubsection*{Rapport}

La classe \texttt{Rapport} encapsule les métadonnées d'un document soumis
à analyse. Elle stocke le nom du fichier, le format (\texttt{pdf},
\texttt{docx}, \texttt{xlsx}), le chemin de stockage sur disque (chaîne
de type \texttt{reports/\{uuid\}/\{nom\}}), et le statut du cycle de vie
(\textit{en\_attente} → \textit{en\_cours} → \textit{terminé} ou
\textit{erreur}). La suppression d'un rapport entraîne la suppression
en cascade de toutes ses analyses associées
(\texttt{cascade="all, delete-orphan"}).

\subsubsection*{Analyse}

La classe \texttt{Analyse} constitue le cœur du modèle. Elle est liée à
un rapport (many-to-one) et à un utilisateur déclencheur. Elle porte le
résumé exécutif (texte libre), le score de maturité GRC (flottant entre
0 et 100), et le statut de traitement. Elle agrège en composition les
risques, les résultats de conformité et les recommandations.

\subsubsection*{Risque}

La classe \texttt{Risque} représente une menace, vulnérabilité ou écart
identifié par Mistral dans le document. Ses attributs clés sont :
\begin{itemize}
  \item \texttt{libelle} — intitulé du risque (500 caractères max)
  \item \texttt{categorie} — taxonomie en cinq valeurs : CYBER,
    OPERATIONNEL, LEGAL, FINANCIER, RH
  \item \texttt{probabilite} et \texttt{impact} — entiers de 1 à 5
  \item \texttt{score\_risque} — produit probabilité × impact (flottant)
  \item \texttt{severite} — calculée automatiquement selon les seuils
    définis dans \texttt{compute\_severity()}
  \item \texttt{section\_source} — référence à l'identifiant ou à la
    section du document source (ex. : \og{}Finding \#3\fg{},
    \og{}Section A.12\fg{})
\end{itemize}

\subsubsection*{ResultatConformite}

La classe \texttt{ResultatConformite} enregistre l'évaluation d'un domaine
de conformité pour un référentiel donné. Ses attributs principaux sont
le référentiel (\texttt{ISO27001}, \texttt{RGPD}, \texttt{LOI0908}), le
nom du domaine évalué, le statut (\texttt{CONFORME} / \texttt{NON\_CONFORME}
/ \texttt{PARTIEL}), l'écart constaté (texte libre) et le taux de
conformité (flottant). Une analyse peut produire jusqu'à 37 enregistrements
de conformité (14 domaines ISO + 12 RGPD + 11 Loi 09-08).

\subsubsection*{Recommandation}

La classe \texttt{Recommandation} modélise une action corrective générée
par Mistral en réponse à un risque. Elle est liée optionnellement à un
risque (\texttt{risque\_id} nullable), permettant la traçabilité
risque-recommandation. Ses attributs incluent : la priorité (CRITIQUE /
HAUTE / MOYENNE / FAIBLE), le type d'action (QUICK\_WIN ou LONG\_TERME),
l'effort estimé (FAIBLE / MOYEN / ÉLEVÉ) et le statut de traitement
(À\_FAIRE / EN\_COURS / CLÔTURÉ).

\subsubsection*{MessageChat}

La classe \texttt{MessageChat} (présente dans le modèle conceptuel) stocke
l'historique des interactions conversationnelles : la question posée par
l'utilisateur, la réponse générée par Mistral et les sources citées
(tableau JSON des extraits ChromaDB).

\subsubsection*{ReferentielGRC}

La classe \texttt{ReferentielGRC} représente la configuration des cadres
normatifs : ISO 27001, RGPD et Loi 09-08. Elle stocke le code du
référentiel, sa version, et la liste de ses domaines et exigences au
format JSON.

\subsubsection*{PisteAudit}

La classe \texttt{PisteAudit} enregistre les actions réalisées par les
utilisateurs (upload, lancement d'analyse, modification de statut),
permettant à l'administrateur de tracer les opérations sensibles via le
type d'entité, l'identifiant de l'entité concernée et les détails en JSON.

\subsection{Relations entre classes}

Le tableau suivant résume les relations de cardinalité du modèle :

\begin{table}[htbp]
  \centering
  \caption{Relations entre les classes du domaine}
  \begin{tabular}{p{3.2cm}p{3.8cm}p{3.2cm}p{4.8cm}}
    \toprule
    \textbf{Classe source} & \textbf{Relation} & \textbf{Classe cible} & \textbf{Description} \\
    \midrule
    Utilisateur & 1..* (uploade) & Rapport & Un utilisateur peut soumettre plusieurs rapports \\
    Utilisateur & 1..* (lance) & Analyse & Un utilisateur peut déclencher plusieurs analyses \\
    Rapport & 0..1 (fait l'objet de) & Analyse & Un rapport peut être analysé une fois à la fois \\
    Analyse & 0..* (contient) & Risque & Une analyse produit N risques \\
    Analyse & 0..* (produit) & ResultatConformite & Une analyse produit jusqu'à 37 résultats de conformité \\
    Analyse & 0..* (génère) & Recommandation & Une analyse génère une recommandation par risque \\
    Analyse & 0..* (contexte) & MessageChat & Une analyse peut faire l'objet de N questions RAG \\
    Risque & 0..* (justifie) & Recommandation & Un risque peut être lié à plusieurs recommandations \\
    ResultatConformite & basé sur & ReferentielGRC & Les résultats référencent un cadre normatif \\
    Utilisateur & 0..* (génère) & PisteAudit & Les actions utilisateur alimentent la piste d'audit \\
    \bottomrule
  \end{tabular}
\end{table}

% ============================================================
\section{Diagramme de Séquence}
% ============================================================

\subsection{Séquence principale : Analyse d'un rapport GRC}

Le diagramme de séquence (figure~\ref{fig:sequence_diagram}) décrit les
interactions chronologiques entre les composants lors de l'opération la
plus complexe du système : le lancement et l'exécution complète d'une
analyse IA.

\begin{figure}[htbp]
  \centering
  \includegraphics[width=\textwidth]{figures/sequence_diagram.png}
  \caption{Diagramme de séquence — Analyse d'un rapport GRC}
  \label{fig:sequence_diagram}
\end{figure}

\subsection{Description détaillée des interactions}

Le flux se décompose en deux grandes phases, séparées par une frontière
asynchrone matérialisée par la réponse HTTP 202.

\subsubsection*{Phase synchrone : Upload et déclenchement}

\begin{enumerate}
  \item L'analyste soumet son fichier depuis le frontend React via
    \texttt{POST /api/rapports/upload}. Axios joint le token JWT dans
    l'en-tête de la requête.
  \item Le backend FastAPI valide le format du fichier et délègue au
    \texttt{report\_service} qui sauvegarde le fichier sur disque et
    crée l'enregistrement \texttt{Rapport} en base de données
    (statut : \textit{en\_attente}).
  \item Le frontend reçoit le \texttt{rapport\_id} et affiche une
    confirmation d'upload.
  \item L'analyste déclenche l'analyse via
    \texttt{POST /api/analyses/lancer/\{rapport\_id\}}.
  \item Le backend crée immédiatement l'enregistrement \texttt{Analyse}
    (statut : \textit{en\_cours}), enregistre la tâche de fond via
    \texttt{BackgroundTasks}, et retourne HTTP 202 avec
    l'\texttt{analyse\_id}.
\end{enumerate}

\subsubsection*{Phase asynchrone : Pipeline IA (8 appels Mistral)}

Le service \texttt{analysis\_service.run\_analysis()} s'exécute dans
une session SQLAlchemy dédiée. Il suit le pipeline suivant :

\begin{enumerate}
  \item \textbf{Extraction du texte} — le fichier est lu depuis le disque
    et parsé selon son format. Le texte brut est découpé en chunks.

  \item \textbf{Phase 1 — Parallélisation via \texttt{asyncio.gather}} :
    \begin{itemize}
      \item Indexation ChromaDB (embedding des chunks en mémoire)
      \item Prompt 1 → Mistral : extraction exhaustive des risques
        (JSON structuré)
      \item Prompt 2 → Mistral : scoring probabilité × impact des risques
      \item Prompt 3 → Mistral : vérification conformité ISO 27001
        (14 domaines)
      \item Prompt 4 → Mistral : vérification conformité RGPD (12 domaines)
      \item Prompt 5 → Mistral : vérification conformité Loi 09-08
        (11 domaines)
    \end{itemize}
    Les résultats sont validés (sévérités, catégories, doublons) et
    persistés en base de données.

  \item \textbf{Gap Analysis consolidée} — à partir des taux de conformité
    calculés, le score de maturité global est déterminé (moyenne des
    trois référentiels).

  \item \textbf{Phase 2 — Parallélisation via \texttt{asyncio.gather}} :
    \begin{itemize}
      \item Prompt 7 → Mistral : génération des recommandations priorisées
        (une par risque identifié)
      \item Prompt 8 → Mistral : génération du résumé exécutif
        (350-450 mots, structuré en 6 sections)
    \end{itemize}

  \item \textbf{Sauvegarde finale} — l'enregistrement \texttt{Analyse}
    est mis à jour avec le résumé, le score de maturité et le statut
    \textit{terminé}.
\end{enumerate}

\subsubsection*{Phase de consultation}

Le frontend React interroge \texttt{GET /api/analyses/\{id\}} toutes les
4 secondes (polling) jusqu'à obtenir le statut \textit{terminé}.
À ce moment, trois requêtes parallèles (\texttt{Promise.all}) récupèrent
les risques, les résultats de conformité et les recommandations, qui sont
affichés dans les onglets correspondants de la page d'analyse.

% ============================================================
\section{Diagramme d'Activité}
% ============================================================

\subsection{Vue d'ensemble du flux de traitement}

Le diagramme d'activité (figure~\ref{fig:activity_diagram}) modélise le
flux de contrôle complet du système, de l'upload du rapport jusqu'à la
consultation des résultats et l'export des livrables. Il met en évidence
les deux zones de parallélisme (swimlanes horizontaux) correspondant aux
deux phases du pipeline IA.

\begin{figure}[htbp]
  \centering
  \includegraphics[width=0.85\textwidth]{figures/activity_diagram.png}
  \caption{Diagramme d'activité — flux de traitement GRC AI Analyzer}
  \label{fig:activity_diagram}
\end{figure}

\subsection{Description des activités}

\subsubsection*{Bloc 1 — Préparation du document}

Les quatre premières activités forment une séquence linéaire :

\begin{enumerate}
  \item \textbf{Upload du rapport} — l'analyste soumet son fichier
    PDF/DOCX/XLSX via l'interface. Le système valide le format,
    sauvegarde le fichier et enregistre les métadonnées.
  \item \textbf{Extraction et nettoyage du texte} — selon le format
    détecté, l'un des trois parsers est appelé (PyMuPDF pour PDF,
    python-docx pour DOCX, openpyxl pour XLSX). Le texte est extrait
    structuré par page, paragraphe ou cellule.
  \item \textbf{Découpage en chunks} — le texte est fragmenté en segments
    de 1 500 caractères avec un chevauchement de 200 caractères assurant
    la continuité sémantique.
  \item \textbf{Génération des embeddings} — chaque chunk est transformé
    en vecteur dense de 384 dimensions par le modèle
    \texttt{paraphrase-multilingual-MiniLM-L12-v2} et indexé dans
    ChromaDB dans une collection nommée
    \texttt{analyse\_\{id\}}.
\end{enumerate}

\subsubsection*{Bloc 2 — Analyse IA parallèle (Phase 1)}

Un nœud de fourche (\textit{fork}) déclenche simultanément quatre
activités parallèles, matérialisant l'utilisation de
\texttt{asyncio.gather} dans le code :

\begin{itemize}
  \item \textbf{Extraction des risques (Mistral)} — appel JSON structuré
    retournant la liste des risques avec libellé, description, catégorie,
    probabilité, impact et référence source.
  \item \textbf{Scoring probabilité × impact} — calcul déterministe du
    score (\textit{score = probabilité × impact}) et classification
    automatique de la sévérité selon les seuils définis.
  \item \textbf{Vérification ISO 27001 (Mistral)} — évaluation des
    14 domaines de l'Annexe A avec statut CONFORME/PARTIEL/NON\_CONFORME
    et taux par domaine.
  \item \textbf{Vérification RGPD (Mistral)} — évaluation des 12 articles
    RGPD pertinents.
  \item \textbf{Vérification Loi 09-08 (Mistral)} — évaluation des
    11 domaines de la loi marocaine de protection des données.
\end{itemize}

Un nœud de jonction (\textit{join}) synchronise les cinq activités
parallèles avant de passer à la consolidation.

\subsubsection*{Bloc 3 — Consolidation et génération (Phase 2)}

\begin{enumerate}
  \item \textbf{Consolidation gap analysis} — les taux de conformité
    sont agrégés, le score de maturité global est calculé.
  \item \textbf{Génération des recommandations priorisées (Mistral)} et
    \textbf{génération du résumé exécutif (Mistral)} s'exécutent en
    parallèle.
  \item \textbf{Sauvegarde en base de données} — tous les objets (risques,
    conformités, recommandations, résumé) sont persistés dans PostgreSQL.
  \item \textbf{Mise à jour du tableau de bord} — le statut de l'analyse
    passe à \textit{terminé}.
\end{enumerate}

\subsubsection*{Décision : Risques critiques détectés ?}

Après la mise à jour du tableau de bord, un nœud de décision vérifie
la présence de risques de sévérité CRITIQUE (score $\geq$ 20). Si oui,
une notification est envoyée à l'administrateur. Dans les deux cas,
l'analyste est redirigé vers la page des résultats.

\subsubsection*{Bloc 4 — Consultation des résultats (parallèle)}

L'analyste dispose de quatre vues accessibles simultanément via les
onglets de l'interface :
\begin{itemize}
  \item \textbf{Matrice des risques} — visualisation Probabilité × Impact
    via \texttt{RiskMatrix} (Recharts).
  \item \textbf{Taux de conformité par référentiel} — radar chart
    via \texttt{ComplianceRadar} (Recharts).
  \item \textbf{Recommandations} — plan d'action priorisé avec gestion
    du cycle de vie.
  \item \textbf{Chat IA} — assistant conversationnel RAG (Mistral +
    ChromaDB).
\end{itemize}

Puis l'analyste peut suivre l'avancement des recommandations au fil
du temps.

\subsubsection*{Décision : Export nécessaire ?}

Si l'analyste souhaite exporter, un nœud de fourche permet de générer
en parallèle le rapport PDF (ReportLab) et/ou le registre Excel
(XlsxWriter), avant de converger vers la fin du flux.

% ============================================================
\section*{Conclusion}
% ============================================================

Ce chapitre a établi les fondements conceptuels de la plateforme GRC AI
Analyzer. L'analyse des besoins a permis d'identifier 24 besoins
fonctionnels répartis en six modules et 13 besoins non fonctionnels
couvrant la performance, la fiabilité, la sécurité et la maintenabilité.
Le diagramme des cas d'utilisation a formalisé les interactions entre les
trois acteurs humains et le modèle Mistral AI. Le diagramme de classes a
modélisé les huit entités persistantes du domaine et leurs relations de
composition et d'association. Le diagramme de séquence a détaillé les
huit appels successifs à l'API Mistral orchestrés en deux phases de
parallélisme. Enfin, le diagramme d'activité a retranscrit le flux complet
du traitement, depuis l'upload jusqu'à l'export des livrables.

Ces artefacts constituent la référence pour le chapitre suivant, qui
présentera l'implémentation concrète de la solution.
